\subsection{Exact Fourier calculation for the retained theta source}
\label{lem:typed-route-fourier-phi-star}

We prove the Fourier self-duality used in the preceding definition, with
the original convention and every coefficient. Put
\(u(x)=e^{-\pi x^2}\). Its integral is one. Indeed, its positive integral
\(I\) satisfies, by Tonelli's theorem and the polar-coordinate Jacobian,
\[
I^2=\int_{\mathbb R^2}e^{-\pi(x^2+y^2)}\,dx\,dy
 =\int_0^{2\pi}\int_0^\infty e^{-\pi r^2}r\,dr\,d\vartheta
 =2\pi\frac1{2\pi}=1,
\]
so \(I=1\).

Differentiation under the Fourier integral is justified by the integrable
majorants \((2\pi|x|)^k u(x)\), for each fixed derivative order \(k\).
Since \(u'=-2\pi xu\), integration by parts gives
\[
\begin{aligned}
\frac{d}{d\xi}\widehat u(\xi)
 &=\int_{\mathbb R}(-2\pi ix)u(x)e^{-2\pi ix\xi}\,dx\\
 &=i\int_{\mathbb R}u'(x)e^{-2\pi ix\xi}\,dx\\
 &=i(2\pi i\xi)\widehat u(\xi)
 =-2\pi\xi\widehat u(\xi).
\end{aligned}
\]
The boundary term is zero because \(u(x)\) tends to zero at both
ends of the real line. The differential equation, together with
\(\widehat u(0)=1\), gives
\[
\widehat u(\xi)=e^{-\pi\xi^2}.
\]
For every Schwartz function \(f\), the same differentiation and integration
by parts prove the typed operator identities
\[
\mathcal F(xf)=\frac{i}{2\pi}\partial_\xi\widehat f,\qquad
\mathcal F(f')=2\pi i\xi\widehat f.
\]
The Gaussian derivatives are
\[
u''(\xi)=(4\pi^2\xi^2-2\pi)u(\xi),
\]
\[
u^{(4)}(\xi)=
(16\pi^4\xi^4-48\pi^3\xi^2+12\pi^2)u(\xi).
\]
Consequently,
\[
\begin{aligned}
\mathcal F(x^2u)(\xi)
 &=-\frac1{4\pi^2}u''(\xi)
   =\left(\frac1{2\pi}-\xi^2\right)u(\xi),\\
\mathcal F(x^4u)(\xi)
 &=\frac1{16\pi^4}u^{(4)}(\xi)
   =\left(\xi^4-\frac3\pi\xi^2+\frac3{4\pi^2}\right)u(\xi).
\end{aligned}
\]
Apply these formulas to the original source
\(\phi_*=(4\pi^2x^4-6\pi x^2)u\):
\[
\begin{aligned}
\widehat{\phi_*}(\xi)
 &=\left[
 4\pi^2\left(\xi^4-\frac3\pi\xi^2+\frac3{4\pi^2}\right)
 -6\pi\left(\frac1{2\pi}-\xi^2\right)\right]u(\xi)\\
 &=(4\pi^2\xi^4-12\pi\xi^2+3-3+6\pi\xi^2)u(\xi)\\
 &=(4\pi^2\xi^4-6\pi\xi^2)e^{-\pi\xi^2}
 =\phi_*(\xi).
\end{aligned}
\]
This establishes the literal equality
\(\widehat{\phi_*}=\phi_*\). The function is even and Schwartz,
\(\phi_*(0)=0\), and
\(\widehat{\phi_*}(0)=\phi_*(0)=0\); hence it belongs to the exact source
space \(V\).

The same calculation also verifies the source operator expression and
the sign of its Fourier intertwiner. Since
\[
Du=2\pi x^2u,\qquad
D^2u=(-4\pi x^2+4\pi^2x^4)u,
\]
we have
\[
D(D-1)u=(4\pi^2x^4-6\pi x^2)u=\phi_*.
\]
For every Schwartz \(f\),
\[
\begin{aligned}
\widehat{Df}(\xi)
 &=-\mathcal F(xf')(\xi)
 =-\frac{i}{2\pi}\partial_\xi(2\pi i\xi\widehat f(\xi))\\
 &=\widehat f(\xi)+\xi\partial_\xi\widehat f(\xi)
 =(1-D_\xi)\widehat f(\xi).
\end{aligned}
\]
Here \(D_\xi=-\xi\partial_\xi\), with the same sign as the original
operator. This proves \(\mathcal F D=(1-D)\mathcal F\) on its stated
Schwartz domain. Moreover, \(D\) preserves \(V\): it preserves evenness
and the Schwartz class, \((D\phi)(0)=0\), and
\[
\int_{\mathbb R}D\phi(x)\,dx
=-[x\phi(x)]_{-\infty}^{\infty}
 +\int_{\mathbb R}\phi(x)\,dx=0
\quad(\phi\in V).
\]
The boundary term vanishes by Schwartz decay. Thus every
\(D^m\phi_*\) and every \((1-D)^m\phi_*\) occurring in the two-leg
calculation remains in the original domain.

The Mellin constant can also be obtained directly from the same
function. For \(\Re s>0\), substitution \(t=\pi x^2\) and
\(\Gamma(z+1)=z\Gamma(z)\) give
\[
\begin{aligned}
M_+\phi_*(s)
&=\int_0^\infty
 (4\pi^2x^4-6\pi x^2)e^{-\pi x^2}x^s\,\frac{dx}{x}\\
&=\pi^{-s/2}
 \left(2\Gamma\left(\frac s2+2\right)
       -3\Gamma\left(\frac s2+1\right)\right)\\
&=\pi^{-s/2}\Gamma\left(\frac s2\right)
 \left[2\frac s2\left(\frac s2+1\right)-3\frac s2\right]\\
&=\frac12s(s-1)\pi^{-s/2}\Gamma\left(\frac s2\right).
\end{aligned}
\]
The recurrence follows by integration by parts in the defining
Gamma integral on this half-plane; its boundary terms vanish
at zero and infinity. For \(\Re s>1\), the full nonzero-integer
theta sum is absolutely integrable term by term, because
\[
\sum_{n\ne0}\int_0^\infty
 |\phi_*(nx)|x^{\Re s}\,\frac{dx}{x}
=2\sum_{n\ge1}n^{-\Re s}
 \int_0^\infty|\phi_*(y)|y^{\Re s}\,\frac{dy}{y}<\infty.
\]
Therefore
\[
\begin{aligned}
\mathcal M\Theta\phi_*(s)
&=2\zeta(s)M_+\phi_*(s)\\
&=s(s-1)\pi^{-s/2}\Gamma(s/2)\zeta(s)
 =2\xi(s)=g(s).
\end{aligned}
\]
This retains both the factor two from the full integer sum and the
factor one half in the source Mellin integral. The source's entire
continuation identifies these same functions away from this initial
half-plane; no division by \(g\) at an arithmetic zero has occurred.

